\section{Argus Life and the Technical Spine}
\label{sec:life}

\begin{designbox}
Turn the report's master spine --- an unknown objective entering a
dense-intelligence runtime, passing an evidence gate, updating durable runtime
state, and meeting the next task from a higher floor --- into named, source-grounded
components. Argus Life keeps the lifetime objective moving; each mission is a
bounded unit with a first-class outcome.
\end{designbox}

Act~I (Section~\ref{sec:episodic}) derived five requirements for a compounding
research system: a persistent objective, role separation, bounded context,
independent evidence, and reusable runtime state. Act~II is the running system that
meets them. This section is the map from the expanding-frontier spine
(Figure~\ref{fig:spine}) to the implementation --- which component owns each stage
and what durable artifact it leaves behind --- and the sections that follow expand
each stage: the four roles and three planes (Section~\ref{sec:rolesplanes}), the
mission lifecycle and durable state (Section~\ref{sec:lifecycle}), the evidence and
review architecture (Section~\ref{sec:evidence}), and the reliability and resource
machinery (Section~\ref{sec:operations}).

\subsection{Life keeps the objective moving; missions stay bounded}

\emph{Argus Life} is the persistent context of a campaign: a durable,
operator-anchored objective and its committed configuration, held on disk and
recovered across restarts and upgrades. A \code{STANDING} objective is persisted
atomically to \code{continuous.json} and is never silently re-decided
(Section~\ref{sec:lifecycle}). Life itself issues no provider work; it exists so a
long campaign has a stable identity. The unit of \emph{work} is a mission: one
bounded Engineer$\leftrightarrow$Reviewer run that always ends in a first-class
terminal or paused status, never an unbounded spin. Life supplies continuity;
missions supply bounded, auditable progress. That division is what lets the system
run for days while every increment stays a discrete, inspectable event.

\subsection{The spine, mapped to owners and durable outputs}

Table~\ref{tab:spine} maps each stage of the master spine to the component that
implements it and the durable artifact it writes. The mapping is deliberately
one-to-one: every conceptual stage has exactly one implemented owner and one place
its result is persisted, so a reader can move from the figure to the code without
guessing.

\begin{table}[t]
\centering
\small
\begin{tabularx}{\linewidth}{@{}p{3.1cm} p{4.7cm} X@{}}
\toprule
\textbf{Master-spine concept} & \textbf{Implemented owner} & \textbf{Durable output} \\
\midrule
Unknown / OOD objective enters &
Manager front door + division (\code{manager/\_core.py}) &
committed vertical + lifetime in \code{continuous.json} \\
Dense-intelligence runtime &
\code{LifeSupervisor} outer loop (\code{life/supervisor/\_core.py}) + \code{SkillLoop} per mission (\code{loop.py}) &
\code{events.jsonl} tape + project journal \\
Four persistent roles &
Manager, Planner, Engineer, Reviewer &
typed role verdicts on the tape \\
Evidence gate &
Reviewer completion authority (\code{reviewer/\_core.py}) &
\code{research.achievement.certified} \\
Runtime evolution &
Reviewer \code{skill\_ops}/\code{wiki\_ops}; Scientist distill-on-miss (\code{skills/}, \code{wiki/}) &
versioned skills + immutable wiki sources \\
Expanded frontier / next task &
\code{LifeSupervisor} next claim; Planner refill (\code{planner/planner.py}) &
refilled \code{backlog.jsonl} inherited by the next mission \\
\bottomrule
\end{tabularx}
\caption{The master spine (Figure~\ref{fig:spine}) mapped to implemented owners and
their durable outputs. Each stage has one owner and one persistence point.}
\label{tab:spine}
\end{table}

Two properties of the mapping matter. First, the \emph{evidence gate} is a single
role --- the Reviewer --- not a hardcoded check (Section~\ref{sec:rolesplanes}); the
authoritative record of a verified achievement is exactly one event on the tape,
\code{research.achievement.certified} (Section~\ref{sec:evidence}). Second,
\emph{runtime evolution} is not one mechanism but several, and the report is
careful not to flatten them into a single arrow.

\subsection{Runtime evolution has four distinct owners}

The spine's runtime-evolution step --- the update
$H(t{+}1)=U(H(t),\text{trajectory},\text{evidence})$ with the model's parameters
held fixed --- is realized by four ownership paths, each carrying a different
authority.

\begin{description}[leftmargin=1.4em, style=nextline]
  \item[Automatic.] Some updates are mechanical. On a matcher miss a Scientist
    distills an immediately usable project-layer skill; after every real Reviewer
    verdict a \emph{source} wiki RoundCard is written verbatim from the verdict; and
    a recency memory preamble is rendered before each mission. None requires further
    judgment (Section~\ref{sec:lifecycle}).
  \item[Model-proposed.] The Engineer \emph{proposes} --- it drafts new skill
    content and, at a turn boundary, emits a curated working-memory handoff. A
    proposal is a candidate, not yet a committed change.
  \item[Reviewer-authored.] Durable knowledge is \emph{committed} by the Reviewer,
    the sole author of \code{skill\_ops} and \code{wiki\_ops} and the validator of
    the canonical \code{checkpoint.json}. A synthesized wiki page is admitted only
    when its evidence spans verify verbatim against a source
    (Section~\ref{sec:evidence}).
  \item[Operator-invoked.] Some evolution is human. An operator edits the machine
    identity, exports or overwrites skills, sets configuration knobs, switches a
    role's backend through the Manager, or leases GPUs --- explicit actions, never
    inferred by the runtime (Section~\ref{sec:operations}).
\end{description}

Keeping these four distinct is a design commitment: the system never presents an
automatic distillation as if a human authorized it, nor a Reviewer-certified skill
as if it were an unchecked model suggestion. Authority over the runtime is always
attributable.
